\documentclass[UTF8,a4paper,11pt]{ctexart}

\usepackage{amsmath,amssymb,amsthm}
\usepackage{geometry}
\usepackage{enumitem}
\usepackage{booktabs,longtable}
\usepackage{xcolor}
\usepackage{hyperref}

\geometry{left=1.85cm,right=1.85cm,top=2.0cm,bottom=2.0cm}
\hypersetup{colorlinks=true,linkcolor=blue!50!black,urlcolor=blue!50!black}
\setlength{\parindent}{0pt}
\setlength{\parskip}{0.35em}
\setlist[enumerate]{leftmargin=1.9em,itemsep=0.35em,topsep=0.35em}
\setlist[itemize]{leftmargin=1.7em,itemsep=0.25em,topsep=0.25em}

\newcommand{\points}[1]{\hfill{\small\textbf{[#1 分]}}}
\newcommand{\OPT}{\mathrm{OPT}}

\title{\textbf{CS216 Algorithm Design and Analysis (H)\\Greedy / DP / Network Flow 专题练习}}
\author{题目在前，参考答案在最后}
\date{2026 年 6 月}

\begin{document}
\maketitle

\textbf{使用建议：}
\begin{itemize}
  \item 每题都按考试大题格式练习：model, algorithm, correctness, complexity。
  \item Greedy 题重点练 exchange / stays-ahead / lower bound。
  \item DP 题重点练状态定义、转移、边界、计算顺序。
  \item Flow 题重点练建图、cut/flow 含义、可行性与整数性。
  \item 答案在最后，请先独立做完再看。
\end{itemize}

\section*{Part I. Greedy 专题}

\subsection*{G1. 最少打点覆盖区间 \points{10}}

给定 $n$ 个闭区间 $[s_i,f_i]$。请选最少数量的实数点，使每个区间至少包含一个被选点。设计 $O(n\log n)$ 算法并证明正确性。

\subsection*{G2. 最多不重叠区间 \points{10}}

给定 $n$ 个区间 $[s_i,f_i)$，选择尽可能多的两两不重叠区间。设计 $O(n\log n)$ 算法并证明正确性。

\subsection*{G3. 最少教室安排课程 \points{10}}

给定 $n$ 门课的时间区间 $[s_i,f_i)$。同一教室不能安排时间重叠的课程。请设计算法求最少需要多少间教室，并给出一种具体安排。

\subsection*{G4. 带截止期的单位作业最大收益 \points{10}}

有 $n$ 个作业，每个作业耗时均为 $1$。作业 $i$ 有截止时间 $d_i$ 和收益 $v_i$，若在时间 $d_i$ 之前完成则获得收益，否则不得收益。每个时间单位最多做一个作业。请设计算法最大化总收益。

\subsection*{G5. 最少加油次数 \points{10}}

一辆车油箱满油可行驶距离 $D$。从起点到终点沿路有加油站，位置为
\[
0=x_0<x_1<\cdots<x_n<x_{n+1}=L.
\]
每次到加油站可加满油。请设计算法判断能否到达终点，并在能到达时最小化加油次数。

\newpage
\section*{Part II. Dynamic Programming 专题}

\subsection*{D1. 加权区间调度 \points{10}}

有 $n$ 个任务，每个任务 $j$ 有开始时间 $s_j$、结束时间 $f_j$ 和收益 $v_j$。选取互不重叠的任务集合，使总收益最大。设计 $O(n\log n)$ 算法。

\subsection*{D2. 0/1 背包 \points{10}}

有 $n$ 个物品，物品 $i$ 重量 $w_i$、价值 $v_i$，背包容量为 $W$。每个物品最多选一次。请设计动态规划算法求最大价值，并说明如何恢复被选物品集合。

\subsection*{D3. 最少硬币数 \points{10}}

给定硬币面值 $a_1,\ldots,a_k$，每种硬币数量无限。给定金额 $M$，求凑出 $M$ 所需的最少硬币数；若不能凑出，输出 impossible。请给出 DP。

\subsection*{D4. 编辑距离 \points{10}}

给定字符串 $X=x_1\cdots x_m$ 和 $Y=y_1\cdots y_n$。允许 insert、delete、replace，每种操作代价为 $1$。请设计 DP 求将 $X$ 变为 $Y$ 的最小代价。

\subsection*{D5. 分段最小二乘 \points{10}}

给定按 $x$ 坐标排序的点 $(x_i,y_i)$，用若干条线段拟合这些点。每开一段有固定费用 $C$，每段的误差为该段内点到最佳拟合直线的平方误差和。设 $e(i,j)$ 表示点 $i,\ldots,j$ 用一条线段拟合的最小误差。请写出 DP。

\newpage
\section*{Part III. Network Flow / Min-Cut 专题}

\subsection*{F1. 二分图最大匹配 \points{10}}

给定二分图 $G=(L\cup R,E)$，请用最大流算法求最大匹配，并证明建图正确性。

\subsection*{F2. 项目选择与依赖 \points{10}}

有若干项目。项目 $i$ 的利润为 $p_i$，可能为正也可能为负。若选择项目 $i$，则必须选择它依赖的所有项目；依赖关系给出为有向边 $(i,j)$，表示选择 $i$ 必须选择 $j$。请设计最大化总利润的算法。

\subsection*{F3. 点不相交路径 \points{10}}

给定有向图 $G=(V,E)$ 和两个点 $s,t$。请设计多项式时间算法求从 $s$ 到 $t$ 的最大点不相交路径条数。除 $s,t$ 外，其余点最多被一条路径使用。

\subsection*{F4. 带上下界的任务分配 \points{10}}

有 $n$ 个工人和 $m$ 个任务。工人 $i$ 最少做 $a_i$ 个任务，最多做 $b_i$ 个任务；任务 $j$ 最少需要 $L_j$ 个工人，最多需要 $U_j$ 个工人。若工人 $i$ 能做任务 $j$，则给出可行边 $(i,j)$。请判断是否存在可行分配。

\subsection*{F5. 软件迁移最大净收益 \points{10}}

有 $n$ 个软件。迁移软件 $i$ 可获得收益 $p_i\ge0$。若一对软件 $(i,j)$ 中恰好迁移一个，则产生损失 $d_{ij}\ge0$。软件 $1$ 不能迁移。请设计一个最大化净收益的算法。

\newpage
\thispagestyle{empty}
\begin{center}
{\Huge \textbf{STOP}}\\[2em]
{\Large 后面是参考答案与解析。}\\[1em]
{\Large 请先独立完成专题练习，再继续往后看。}
\end{center}

\newpage
\section*{参考答案与解析}

\section*{Part I. Greedy 答案}

\subsection*{G1 答案}

按右端点 $f_i$ 从小到大排序。扫描区间，维护最近选点 $x$。若当前区间已包含 $x$，跳过；否则选择点 $f_i$。

正确性：设当前最早结束且未被覆盖的区间为 $[s,f]$。任意可行解必须在该区间中选一个点 $y$。把 $y$ 换成 $f$ 不会减少它对之后区间的覆盖能力，因为之后需要考虑的区间右端点都不小于 $f$。所以存在包含贪心选择的最优解。归纳即可。复杂度 $O(n\log n)$。

\subsection*{G2 答案}

按结束时间 $f_i$ 从小到大排序。每次选择与已选区间兼容且结束最早的区间。

正确性用 stays-ahead：令贪心选出的第 $r$ 个区间为 $g_r$，任意最优解的第 $r$ 个区间为 $o_r$。归纳证明
\[
f(g_r)\le f(o_r).
\]
因此若最优解能选 $k$ 个，贪心也不会在 $k$ 之前停止。复杂度 $O(n\log n)$。

\subsection*{G3 答案}

按开始时间排序。用一个按当前教室结束时间排序的小根堆。扫描课程：
\begin{itemize}
  \item 若堆顶教室结束时间 $\le s_i$，把该课程放入该教室并更新结束时间。
  \item 否则新开一间教室。
\end{itemize}
正确性：若算法开第 $d$ 间教室，说明当前课程开始时已有 $d-1$ 间教室的课程尚未结束，加上当前课程，共有 $d$ 门课同时重叠，所以任何方案至少需要 $d$ 间教室。贪心使用的教室数达到这个下界。复杂度 $O(n\log n)$。

\subsection*{G4 答案}

按截止时间 $d_i$ 从小到大排序。维护一个小根堆，存当前决定按时完成的作业收益。扫描作业 $i$：
\begin{itemize}
  \item 将 $v_i$ 加入堆。
  \item 若堆大小 $>d_i$，弹出收益最小的作业。
\end{itemize}
最后堆中作业为最优选择。

正确性：处理完截止时间不超过 $d$ 的作业后，最多只能按时完成 $d$ 个。若当前选择超过 $d$ 个，保留收益最大的 $d$ 个显然最优；弹出最小收益不会变差。对扫描前缀归纳即可。复杂度 $O(n\log n)$。

\subsection*{G5 答案}

从当前位置出发，总是开到在距离 $D$ 内能到达的最远加油站；若下一个站都到不了，则 impossible；若终点可达则结束。

正确性：设贪心第 $r+1$ 次停在最远可达站 $g_{r+1}$，任意最优解对应停在 $o_{r+1}$。由于 $o_{r+1}$ 从同一位置可达，必有 $g_{r+1}\ge o_{r+1}$。把最优解中的 $o_{r+1}$ 换成 $g_{r+1}$ 不会增加停靠次数，也不会破坏后续可行性。交换归纳得贪心最优。复杂度 $O(n)$。

\newpage
\section*{Part II. DP 答案}

\subsection*{D1 答案}

按结束时间排序。定义
\[
p(j)=\max\{i<j:f_i\le s_j\}.
\]
状态：
\[
M[j]=\text{前 }j\text{ 个任务的最大收益}.
\]
转移：
\[
M[j]=\max\{M[j-1],v_j+M[p(j)]\}.
\]
边界 $M[0]=0$。正确性：最优解要么不选 $j$，要么选 $j$；选 $j$ 时之前只能来自前 $p(j)$ 个任务。排序和二分计算 $p(j)$ 为 $O(n\log n)$，DP 为 $O(n)$。

\subsection*{D2 答案}

状态：
\[
dp[i][w]=\text{只考虑前 }i\text{ 个物品，容量 }w\text{ 的最大价值}.
\]
转移：
\[
dp[i][w]=
\begin{cases}
dp[i-1][w], & w_i>w,\\
\max\{dp[i-1][w],dp[i-1][w-w_i]+v_i\}, & w_i\le w.
\end{cases}
\]
边界 $dp[0][w]=0$。答案 $dp[n][W]$。恢复方案：从 $(n,W)$ 倒推，若 $dp[i][w]=dp[i-1][w]$ 则不选 $i$；否则选 $i$ 并令 $w\leftarrow w-w_i$。复杂度 $O(nW)$。

\subsection*{D3 答案}

状态：
\[
dp[x]=\text{凑出金额 }x\text{ 的最少硬币数}.
\]
边界 $dp[0]=0$，其余初始化为 $+\infty$。转移：
\[
dp[x]=1+\min_{a_i\le x}dp[x-a_i].
\]
按 $x=1,\ldots,M$ 递增计算。若 $dp[M]=+\infty$，输出 impossible。复杂度 $O(kM)$。

\subsection*{D4 答案}

状态：
\[
dp[i][j]=\text{将 }x_1\cdots x_i\text{ 变为 }y_1\cdots y_j\text{ 的最小代价}.
\]
边界：
\[
dp[i][0]=i,\qquad dp[0][j]=j.
\]
转移：
\[
dp[i][j]=\min\begin{cases}
dp[i-1][j]+1, & \text{delete }x_i,\\
dp[i][j-1]+1, & \text{insert }y_j,\\
dp[i-1][j-1]+[x_i\ne y_j], & \text{match/replace}.
\end{cases}
\]
复杂度 $O(mn)$。

\subsection*{D5 答案}

预处理每个 $e(i,j)$。状态：
\[
M[j]=\text{拟合前 }j\text{ 个点的最小总代价}.
\]
边界 $M[0]=0$。最后一段一定是某个后缀 $i,\ldots,j$，所以
\[
M[j]=\min_{1\le i\le j}\{M[i-1]+e(i,j)+C\}.
\]
正确性：枚举最优解最后一段起点 $i$，前面部分必须是前 $i-1$ 个点的最优解。若 $e(i,j)$ 已知，DP 为 $O(n^2)$。

\newpage
\section*{Part III. Network Flow 答案}

\subsection*{F1 答案}

建源点 $s$、汇点 $t$。对每个 $u\in L$ 加边 $s\to u$，容量 $1$；对每条二分图边 $(u,v)$ 加边 $u\to v$，容量 $1$；对每个 $v\in R$ 加边 $v\to t$，容量 $1$。求最大流，所有流量为 $1$ 的 $L\to R$ 边构成匹配。

正确性：容量限制保证每个左点和右点最多匹配一次。任意匹配对应同等大小整数流；任意整数流对应同等大小匹配。整数容量保证存在整数最大流。复杂度为一次最大流。

\subsection*{F2 答案}

这是最大权闭合图。建图：若 $p_i>0$，加边 $s\to i$ 容量 $p_i$；若 $p_i<0$，加边 $i\to t$ 容量 $-p_i$；对依赖边 $(i,j)$ 加边 $i\to j$ 容量 $M$，其中 $M$ 足够大。求最小割，取源侧点集 $S$。

若 $i\in S$ 但依赖的 $j\notin S$，割会切容量 $M$ 的边，不可能最优，所以 $S$ 满足依赖闭合。割容量等于未选正利润损失加已选负利润代价，因此最小割等价于最大总利润。

\subsection*{F3 答案}

拆点。对每个 $v\ne s,t$，拆成 $v_{\mathrm{in}}$ 和 $v_{\mathrm{out}}$，加边 $v_{\mathrm{in}}\to v_{\mathrm{out}}$ 容量 $1$。对 $s,t$ 可设容量无穷或不拆。原边 $(u,v)$ 改为 $u_{\mathrm{out}}\to v_{\mathrm{in}}$，容量无穷；涉及 $s,t$ 时用对应源汇端点。求 $s$ 到 $t$ 最大流。

每个中间点的拆点边容量 $1$，保证最多一条路径使用该点。整数流可分解为点不相交路径。最大流值即最大条数。

\subsection*{F4 答案}

建带下界网络。源点 $s$ 到工人 $u_i$ 加边容量 $[a_i,b_i]$；若工人 $i$ 能做任务 $j$，加边 $u_i\to v_j$ 容量 $[0,1]$；任务 $v_j$ 到汇点 $t$ 加边 $[L_j,U_j]$；再加边 $t\to s$ 容量 $[0,\infty]$ 转成 circulation。

处理下界：每条边先送下界 $\ell_e$，剩余容量 $u_e-\ell_e$。计算
\[
b(v)=\sum_{\text{in}}\ell_e-\sum_{\text{out}}\ell_e.
\]
若 $b(v)>0$ 加 $v\to TT$ 容量 $b(v)$；若 $b(v)<0$ 加 $SS\to v$ 容量 $-b(v)$。若最大流能饱和所有 $SS$ 出边，则可行。流量为 $1$ 的工人到任务边即分配。

\subsection*{F5 答案}

令 $S$ 为迁移的软件集合，且 $1\notin S$。目标：
\[
\max_S\left(\sum_{i\in S}p_i-\sum_{\text{split }(i,j)}d_{ij}\right).
\]
建最小割图：对每个软件加点。加边 $s\to i$ 容量 $p_i$。对每对有关联的软件 $(i,j)$，加无向惩罚边，可用两条有向边 $i\to j$、$j\to i$，容量 $d_{ij}$。为强制软件 $1$ 不迁移，加边 $1\to t$ 容量 $M$，$M$ 足够大。

取最小割源侧为迁移集合。若正收益点不在源侧，会切掉 $s\to i$，表示放弃收益；若一对软件分居两侧，会切惩罚边，表示支付损失；软件 $1$ 若在源侧会切无限边，因此不会被选。最小化割容量等价于最大化净收益。

\end{document}
